\chapter{镜子与反射}\label{ch:g5:mirrors-reflection}

我们已经证明，光从不绕过拐角弯曲——然而每天早晨，
你的影像都会从浴室的墙上望着你；去年我们承诺过，
还要教一个看见拐角的办法。这办法不是把路弯过去，
而是把路\emph{折}起来。一切都折叠在
一个非凡的表面上：\omterm{def:g5:mirrors-reflection:reflection}{镜子}。

\section{弹回}

\begin{definition}[反射]\label{def:g5:mirrors-reflection:reflection}
光射到某个表面上，既不穿过去，
也不在那里消失，而是弹回来，我们称之为
\emph{反射}\index{反射}：光被反射。
\emph{镜子}\index{镜子}是一种非常光滑的表面，它能以十分有序的
方式反射光——例如玻璃后面的抛光金属、
静止的水面、一把闪亮的勺子。
\end{definition}

\begin{proposition}[弹回规律]\label{prop:g5:mirrors-reflection:bounce}
一束\omterm{def:g4:light-straight-lines:ray}{光线}在\omterm{def:g5:mirrors-reflection:reflection}{镜子}上弹回，就像一个完美的球撞到墙上后弹回：
它离开时的陡峭程度与来时完全相同，
只是在另一侧。直射进来，直射回去；斜射进来，
以相同的斜度斜射出去。入射角与
\omterm{def:g5:mirrors-reflection:reflection}{反射}角总是相等。
\end{proposition}

\begin{omfigure}
\begin{tikzpicture}[scale=0.85]
  \draw[very thick] (0.5,0) -- (10.5,0);
  \foreach \x in {1.0,1.75,...,10.0} {
    \draw[thin] (\x,0) -- (\x-0.3,-0.3);
  }
  \draw[-Stealth, very thick, omThm] (2.0,3.0) -- (5.5,0.05);
  \draw[-Stealth, very thick, omThm] (5.5,0.05) -- (9.0,3.0);
  \draw[dashed] (5.5,0) -- (5.5,3.2);
  \draw[omExo, thick] (4.42,0.95) arc (140:90:1.45);
  \draw[omExo, thick] (6.58,0.95) arc (40:90:1.45);
  \node[font=\small] at (4.7,1.75) {入射};
  \node[font=\small] at (6.35,1.75) {反射};
  \node[below, font=\small] at (5.5,-0.45)
    {法线两侧的角相等};
\end{tikzpicture}

{\small 弹回规律：\omterm{def:g5:mirrors-reflection:reflection}{反射}\omterm{def:g4:light-straight-lines:ray}{光线}与入射\omterm{def:g4:light-straight-lines:ray}{光线}
关于法线对称——法线两侧的角相等。}
\end{omfigure}

\begin{example}[瞄准弹回的光]\label{ex:g5:mirrors-reflection:aiming}
由于弹回十分有序，它可以\emph{瞄准}：
在阳光下拿一面小\omterm{def:g5:mirrors-reflection:reflection}{镜子}，你就能把一个光斑投到任何你选定的墙上——
\omterm{def:g5:mirrors-reflection:reflection}{镜子}一倾斜，光斑便随之移动。遇险的水手和徒步者
正是这样用的：把\omterm{def:g5:mirrors-reflection:reflection}{镜子}对准远处的船只或直升机闪烁，
它就成了一个在许多千米之外都能看见的信号，
而且完全依靠太阳提供\omterm{def:g5:everyday-energy:energy}{能量}。
\end{example}

\begin{example}[为什么墙不是镜子]\label{ex:g5:mirrors-reflection:walls}
白墙也能\omterm{def:g5:mirrors-reflection:reflection}{反射}大量的光，否则房间会很昏暗，
那么它为什么不显示像呢？近看，它的表面
是一片细小凸起的地形：每条入射\omterm{def:g4:light-straight-lines:ray}{光线}都
在它自己那一小块\omterm{def:g5:simple-machines:ramp}{斜面}上遵守弹回规律，整齐的\omterm{def:g4:light-straight-lines:ray}{光线}大军
向各个方向散开。\omterm{def:g5:mirrors-reflection:reflection}{镜子}的光滑使\omterm{def:g5:simple-machines:ramp}{斜面}保持平整，
\omterm{def:g4:light-straight-lines:ray}{光线}大军保持步调一致，画面便能在\omterm{def:g5:mirrors-reflection:reflection}{反射}中保存下来。
决定能否成像的是有序，而不只是光亮。
\end{example}

\begin{omfigure}
\includegraphics[width=0.45\linewidth]{images/book1/ai/g5-mirrors-lake.jpg}

{\small 静止的湖水就像平躺着的\omterm{def:g5:mirrors-reflection:reflection}{镜子}：每条\omterm{def:g4:light-straight-lines:ray}{光线}都按同一规律\omterm{def:g5:mirrors-reflection:reflection}{反射}，于是山的影像再次出现，上下颠倒，分毫不差。}
\end{omfigure}

\section{玻璃后面的人}

\begin{proposition}[成像规律]\label{prop:g5:mirrors-reflection:image}
平面镜会把每个物体显示为\emph{像}\index{像}，
这个像位于玻璃\emph{后面}，与物体在前方的距离相等，
大小相同，一一对应。你后退一步，
你的像也会后退一步；镜中世界
是我们这个世界忠实的、折叠过来的副本。
\end{proposition}

\begin{omfigure}
\begin{tikzpicture}[scale=0.85]
  \draw[very thick] (5.5,0) -- (5.5,3.6);
  \fill[omThm] (2.5,2.6) circle (0.16);
  \draw[thick] (2.5,0.8) -- (2.5,2.4);
  \node[below, font=\small] at (2.5,0.5) {蜡烛，前方\qty{1}{m}};
  \fill[omThm, opacity=0.35] (8.5,2.6) circle (0.16);
  \draw[thick, omExo, dashed] (8.5,0.8) -- (8.5,2.4);
  \node[below, font=\small] at (8.5,0.5) {像，后方\qty{1}{m}};
  \draw[omThm, thick] (2.5,2.6) -- (5.5,1.9);
  \draw[-Stealth, omThm, thick] (5.5,1.9) -- (3.4,0.35);
  \draw[omExo, dashed] (5.5,1.9) -- (8.5,2.6);
  \node[font=\small] at (3.2,0.15) {眼睛};
\end{tikzpicture}

{\small 成像规律：眼睛接收到的是\omterm{def:g5:mirrors-reflection:reflection}{反射}\omterm{def:g4:light-straight-lines:ray}{光线}，但会
沿直线反向追溯——看到的仿佛是一支站在
玻璃后面的蜡烛，与真实蜡烛在前方的距离相等。}
\end{omfigure}

\begin{example}[为什么像在后面]\label{ex:g5:mirrors-reflection:why}
你的眼睛很单纯：它总以为每条\omterm{def:g4:light-straight-lines:ray}{光线}都沿直线传播。
\omterm{def:g5:mirrors-reflection:reflection}{反射}\omterm{def:g4:light-straight-lines:ray}{光线}确实是从蜡烛经\omterm{def:g5:mirrors-reflection:reflection}{镜子}而来，
但眼睛会把它\emph{穿过玻璃向后}延长，
直到\omterm{def:g4:light-straight-lines:ray}{光线}仿佛起始的地方。所有\omterm{def:g5:mirrors-reflection:reflection}{反射}\omterm{def:g4:light-straight-lines:ray}{光线}
都指向同一个虚构的起点：一支蜡烛发出的\omterm{def:g4:light-straight-lines:ray}{光线}，
汇成一支清晰的虚像蜡烛。当然，墙后面什么也没有——
像是光最有说服力的错觉。
\end{example}

\begin{example}[交换的世界]\label{ex:g5:mirrors-reflection:swap}
把一页写着字的纸放到\omterm{def:g5:mirrors-reflection:reflection}{镜子}前：字母会反着回来，
左右互换。这就是为什么
救护车车头上的 AMBULANCE 要画成镜像字样：
前方的司机看向后视镜，会看到它转回正常字样，
——于是让到一旁。两面\omterm{def:g5:mirrors-reflection:reflection}{镜子}可以抵消这种互换：
两次\omterm{def:g5:mirrors-reflection:reflection}{反射}把互换再翻转一次。
\end{example}

\section{有意折叠的路线}

\begin{method}[制作潜望镜]\label{met:g5:mirrors-reflection:periscope}
兑现前面承诺的拐角观察器：
\begin{enumerate}
  \item 取一个又高又窄的盒子；在一个面的靠近顶部处开一个窗口，
        在相对的那个面的靠近底部处再开一个窗口；
  \item 在每个窗口内侧后方各固定一面小\omterm{def:g5:mirrors-reflection:reflection}{镜子}，
        每面\omterm{def:g5:mirrors-reflection:reflection}{镜子}都倾斜成直角的一半，
        两面\omterm{def:g5:mirrors-reflection:reflection}{镜子}的倾斜方向互相平行；
  \item 让上方的窗口对准墙外（或拐角处），
        从下方的窗口向里看。
\end{enumerate}
光的路线会折两次：先沿两面\omterm{def:g5:mirrors-reflection:reflection}{镜子}之间的盒子向下，
再射入你的眼睛。潜艇上的船员正是用这种装置观察海面——
只是把它做得又长又密封。
\end{method}

\begin{omfigure}
\begin{tikzpicture}[scale=0.85]
  \draw[thick] (4.0,0) rectangle (5.6,4.4);
  \draw[very thick, omDef] (4.15,3.4) -- (5.45,4.25);
  \draw[very thick, omDef] (4.15,0.15) -- (5.45,1.0);
  \draw[-Stealth, omThm, very thick] (1.0,3.83) -- (4.8,3.83);
  \draw[-Stealth, omThm, very thick] (4.8,3.83) -- (4.8,0.58);
  \draw[-Stealth, omThm, very thick] (4.8,0.58) -- (8.4,0.58);
  \node[left, font=\small] at (0.9,3.83) {景物发出的光};
  \node[right, font=\small] at (8.5,0.58) {射入你的眼睛};
  \node[right, font=\small, omDef] at (5.7,4.0) {倾斜的镜子};
  \node[right, font=\small, omDef] at (5.7,1.2) {倾斜的镜子};
\end{tikzpicture}

{\small 潜望镜内部：两面平行的倾斜\omterm{def:g5:mirrors-reflection:reflection}{镜子}使光的路线
折两次——先越过墙，再射入你的眼睛。}
\end{omfigure}

\begin{remark}[哈哈镜里的勺子]\label{rem:g5:mirrors-reflection:curved}
向一把光亮的勺子的凹面里看：你的像是倒着的。
把勺子翻过来：你的像又是正的，但被拉长了。弯曲的
镜面在每一点上仍然遵守\omterm{def:g5:mirrors-reflection:reflection}{反射}规律——只是它们的
斜率随曲面改变，因此有序的\omterm{def:g5:mirrors-reflection:reflection}{反射}\omterm{def:g4:light-straight-lines:ray}{光线}会聚或发散，
像会缩小、放大或倒转。弯曲的\omterm{def:g5:mirrors-reflection:reflection}{镜子}
也能聚光：弯曲的碟形镜可以把阳光聚起来煮熟一锅食物——
还记得\omterm{def:g5:everyday-energy:energy}{能量}那一章里的太阳灶吗？它们精确的
规律要留到以后的年级再学；这把勺子只是它们的一个真实预告。
\end{remark}

\section{练习}

\begin{exercise}[$\star$]\label{exo:g5:mirrors-reflection:1}
叙述\omterm{def:g5:mirrors-reflection:reflection}{反射}定律。一条正对镜面射来的\omterm{def:g4:light-straight-lines:ray}{光线}会怎样？
\end{exercise}

\begin{exercise}[$\star$]\label{exo:g5:mirrors-reflection:2}
白墙会\omterm{def:g5:mirrors-reflection:reflection}{反射}大量\omterm{def:g4:light-straight-lines:ray}{光线}，但为什么它不能成像，
而光滑的\omterm{def:g5:mirrors-reflection:reflection}{镜子}却能？
\end{exercise}

\begin{exercise}[$\star$]\label{exo:g5:mirrors-reflection:3}
你站在平面镜前 \qty{2}{m} 处。你的像在哪里，
它离你有多远？
\end{exercise}

\begin{exercise}[$\star$]\label{exo:g5:mirrors-reflection:4}
你放慢脚步走向镜中的像。像会怎样运动，
你与像之间的距离会以多快的速度缩小？
\end{exercise}

\begin{exercise}[$\star$]\label{exo:g5:mirrors-reflection:5}
为什么救护车车头上的 AMBULANCE 要画成反的？
它对谁来说是正的，又要通过什么来看？
\end{exercise}

\begin{exercise}[$\star$]\label{exo:g5:mirrors-reflection:6}
在潜望镜中，光的路线折了几次，
各次折叠起什么作用？\omterm{def:g5:mirrors-reflection:reflection}{镜子}应以怎样的倾角放置？
\end{exercise}

\begin{exercise}[$\star$]\label{exo:g5:mirrors-reflection:7}
一位带着袖珍\omterm{def:g5:mirrors-reflection:reflection}{镜子}的徒步者，如何向远处的直升机发出信号？
这里用到了本章的哪两个概念——说出决定闪光方向的规律，
以及为闪光提供动力的\omterm{def:g5:everyday-energy:store}{能量储存}。
\end{exercise}

\begin{exercise}[$\star$]\label{exo:g5:mirrors-reflection:8}
静止的水面映出群山；风吹皱湖面，画面
碎成闪烁的光斑。请结合
\cref{ex:g5:mirrors-reflection:walls}解释这两种状态。
\end{exercise}

\begin{exercise}[$\star\star$]\label{exo:g5:mirrors-reflection:9}
眼睛会把\omterm{def:g4:light-straight-lines:ray}{光线}“向后延长”。利用这一点解释，为什么你的
像站在\omterm{def:g5:mirrors-reflection:reflection}{镜子}所挂的墙的\emph{后面}——位于
一个肯定没有人站着的地方。从你鼻子\omterm{def:g5:mirrors-reflection:reflection}{反射}出的
所有\omterm{def:g4:light-straight-lines:ray}{光线}在什么问题上取得一致？
\end{exercise}

\begin{exercise}[$\star\star$]\label{exo:g5:mirrors-reflection:10}
一家商店希望门口的顾客能越过看不见的拐角，
看到过道里的情况。一面平面镜须装在两条
走廊的交汇处。它应当怎样倾斜，为什么同样的
安装能让过道看到门口？（光路
可以双向通行。）
\end{exercise}

\begin{exercise}[$\star\star$]\label{exo:g5:mirrors-reflection:11}
镜中的字会反过来；两次\omterm{def:g5:mirrors-reflection:reflection}{反射}可以把
“翻转”再翻转一次。请解释，为什么在
\cref{met:g5:mirrors-reflection:periscope}中的潜望镜，其两面\omterm{def:g5:mirrors-reflection:reflection}{镜子}
使所见景物方向正确，而不是镜像反转。
\end{exercise}
